\section*{Chapter \ref{ch:s1:trend-following} --- Trend Following}

\begin{solution}{exo:s1:trend-following:1}
$\lambda = 60/61 = 0.9836$; the weights halve after $\ln 0.5/\ln \lambda = 42$ days.
\end{solution}

\begin{solution}{exo:s1:trend-following:2}
$0.4/(40 \times 0.25) = 0.04$, 4\% of capital in the commodity; $0.4/(40 \times 0.06) = 0.167$, 16.7\% in the bond future. Each contributes 1\% of
volatility.
\end{solution}

\begin{solution}{exo:s1:trend-following:3}
Over each up-and-down cycle the total move is zero, so the first term, $\tfrac12(\sum r)^2$, is zero, and the book loses half the realised variance: it
buys after the rise and sells after the fall. Choppy markets are the trend follower's cost, as time decay is a straddle holder's.
\end{solution}

\begin{solution}{exo:s1:trend-following:4}
$2.63 \times 0.09 + 0.03 \times (-0.30) - 0.01 = 0.218$: about 22\% in the quarter, if the fall is gradual like the synthetic crashes.
\end{solution}

\begin{solution}{exo:s1:trend-following:5}
The one-month rule turns its gross over 49 times a year and the 250-day breakout 3 times; at a few basis points per unit traded, costs grow with
turnover while the fast rule's gross edge is small. The breakout's Sharpe ratio drops only from 0.80 to 0.78.
\end{solution}

\begin{solution}{exo:s1:trend-following:6}
A decade's Sharpe ratio has a standard error of about $1/\sqrt{10} = 0.32$, so most decades fall between $0.4 - 0.32 = 0.08$ and $0.4 + 0.32 = 0.72$. A
decade at 0.1 is within that range: on its own it is weak evidence that the edge has gone. Look for a mechanism (correlations, crowding, costs) before
concluding.
\end{solution}

\begin{solution}{exo:s1:trend-following:7}
Twelve months is still best (0.42), then six (0.39), then three (0.13). A three-month lookback would follow the faster drift, but it measures it with
much more noise: the past return over $L$ days has a drift part growing with $L$ and a noise part growing with $\sqrt{L}$. At this drift strength the
noise dominates short lookbacks, so the best lookback stays long even though the trends are faster; all Sharpe ratios fall.
\end{solution}

\begin{solution}{exo:s1:trend-following:8}
The best of 200 correlated trials is selected on noise: its expected Sharpe ratio is inflated by the maximum of many estimates, each with a standard
error of about 0.2 over 25 years. Deflate it (Book~7, chapter~20), or blend the rules and judge the blend; out of sample, expect much less.
\end{solution}

\begin{solution}{pb:s1:trend-following:1}
\begin{enumerate}
\item Trading each market in the direction of its own recent move; the sign of the past $L$-day return; a position at new $L$-day highs or lows, held
until the opposite.
\item Past return (equal weights over $L$ days); moving-average crossover (a triangle of weights, more on recent days); breakout (only new extremes).
\item $\sum_t w_t r_t = \tfrac12(\sum_t r_t)^2 - \tfrac12\sum_t r_t^2$, from expanding the square.
\item The square of persistent drifts and positive autocorrelation of returns.
\item AR(1) drifts with a half-life of a year and a standard deviation of half the volatility; carries of which half is earned; volatility clustering by
class; three hundred-day crashes.
\item Signal times 40\% over 40 markets over an EWMA volatility forecast; 2, 1, 2 and 4 basis points per unit traded by class.
\item Past return $-0.00$, 0.26, 0.66; crossovers 0.10, 0.43, 0.72; breakouts 0.17, 0.34, 0.78.
\item The planted drift is slow; blending costs the difference to the best rule (0.55 for nine rules against 0.78) and buys insurance against
choosing the wrong speed.
\item Sizing to a forecast volatility; 0.40 scaled against 0.31 at equal notional.
\item It raises volatility from 7.7\% to 10.2\% and lowers the Sharpe ratio from 0.40 to 0.34: it levers up when markets disagree and the edge is small.
\item Positive returns in prolonged stock declines; $+2.3\%$, $+21.1\%$ and $+39.0\%$ while equities lost 6.8\%, 24.3\% and 35.1\%.
\item Quarterly returns against equities are U-shaped, curvature 2.63; the worst equity quarter was the book's best.
\item \textbf{Named result.} In the three synthetic crashes the trend book at 10\% volatility returned $+2.3\%$, $+21.1\%$ and $+39.0\%$ against equity
losses of 6.8\%, 24.3\% and 35.1\%; the Sharpe ratio after costs rises with speed from $-0.00$ (one month) to 0.66 (twelve months) for past returns,
0.10 to 0.72 for crossovers and 0.17 to 0.78 for breakouts.
\item 9.0\% a year at a Sharpe ratio of 0.22, 1986--2024; with a standard error of about 0.16, not much.
\item Persistence for one to twelve months in 58 markets and substantial abnormal returns; positive returns in every decade since 1880 and good
performance in eight of ten crises.
\item The level of pairwise correlations between markets.
\item A standard error of about 0.32.
\item A sudden crash, such as 19 October 1987, faster than the signals' lookback.
\item Chapter~5 ranks stocks against each other and is market-neutral; here each market is traded on its own past, and the book can be net long or short.
\item Large moves, in either direction, paid for with losses in choppy markets.
\end{enumerate}
\end{solution}

\begin{solution}{iq:s1:trend-following:1}
\omterm{def:s1:trend-following:trend}{Time-series momentum} trades each asset on its own past return and can be net long or short the whole market; \omterm{def:s1:momentum:cross}{cross-sectional momentum} ranks assets
against each other, long winners and short losers, and is neutral by construction.
\end{solution}

\begin{solution}{iq:s1:trend-following:2}
So that each market contributes similar risk: without it the most volatile markets dominate and the book is not diversified; and so that the book's risk
is steadier through volatility regimes.
\end{solution}

\begin{solution}{iq:s1:trend-following:3}
Compare the loss with the distribution of two-year returns implied by the book's expected Sharpe ratio; check whether the mechanism has changed (correlations,
crowding, costs, market access); check the implementation (rolls, costs, sizing). Keep it if the loss is within the noise and the mechanism is intact.
\end{solution}

\begin{solution}{iq:s1:trend-following:4}
Contract rolls and continuous series; expiry calendars; margin instead of cash; contract multipliers and tick sizes; returns on notional, not on
capital; different trading hours and holidays by market.
\end{solution}

\begin{solution}{iq:s1:trend-following:5}
It has hedged prolonged declines, because its signals turn short as the decline lasts; it does not hedge sudden crashes, reversals after a long rise, or
choppy markets in which both stocks and the trend book lose.
\end{solution}

\begin{solution}{iq:s1:trend-following:6}
With independent returns, $E[w_t r_t] = \sum_{k} E[r_{t-k}]E[r_t] = L\mu^2$. For small $\mu/\sigma$, $\operatorname{Var}(w_t r_t) \approx E[w_t^2]E[r_t^2]
\approx L\sigma^2 \cdot \sigma^2$, so the daily Sharpe ratio is $L\mu^2/(\sqrt{L}\sigma^2) = \sqrt{L}(\mu/\sigma)^2$, times $\sqrt{252}$ a year: it grows
with the square root of the lookback and the square of the drift's Sharpe ratio.
\end{solution}
